% Originally: content/week-10.tex, \section{A short introduction to geodesics} (Corsin Week 10)


\section{A short introduction to geodesics}
\label{sec:geodesics}

% Checked 2026-08-10 against lec_notes.pdf (239 pp, 27 July 2026): "geodesic" occurs exactly
% once in the whole file, in Example 9.14 on printed p. 5, and the surrounding pages were
% rendered and read rather than grepped. There is no geodesics section, no length functional
% and no variational argument anywhere in the lecture notes.
\begin{ainote}[This chapter goes beyond the course]
\label{ainote:geodesics_beyond_the_course}
Two warnings before starting, and they point the same way.

Corsin marks this section \qt{not exam relevant}. And the official lecture notes have no
counterpart to it at all: the word \newterm{geodesic} appears in them exactly once, in
Example 9.14 on printed p.\ 5, where the great-circle distance
$d_1(x,y) = R\arccos\big(x\cdot y/R^2\big)$ is offered as a more natural metric on the sphere
than the Euclidean one, with no indication of why that arc is the shortest. Nothing in the
course computes a geodesic, and the calculus of variations does not appear in it.

So this chapter is not lecture material, and nothing in it can be assumed known. It is here
because the calculation is short and because the idea behind it, treating length as a functional
and applying optimization to it, is worth meeting once; and it is here in particular because it
answers the question Example 9.14 raises and then drops.
\end{ainote}

We will try to find the shortest path (i.e., a geodesic) from $0 \in \mathbb{R}^n$ to $x \in \mathbb{R}^n$. This
is, of course, the straight line, but the essentially same calculation can be used to find the
shortest path between points on a sphere or any other submanifold.

\textbf{Idea.} Apply \qt{optimization} to the map
\[L : \mathcal A \to [0,\infty), \qquad c \mapsto L(c) = \int_0^1 \lvert c'(t)\rvert\,dt,
  \qquad \mathcal A := \{\gamma \in C^1([0,1],\mathbb{R}^n) : \gamma(0)=0,\ \gamma(1)=x\}.\]

\begin{lemma}[A few differentiation identities]
\label{lem:inner_product_derivative_identities}
Let $\alpha,\beta : (0,1) \to \mathbb{R}^n$ be $C^1$. Then
\[\frac{d}{dt}\langle\alpha(t),\beta(t)\rangle = \langle\alpha'(t),\beta(t)\rangle
  + \langle\alpha(t),\beta'(t)\rangle, \qquad
  \frac{d}{dt}\lvert\alpha(t)\rvert = \frac{1}{\lvert\alpha(t)\rvert}\langle\alpha(t),\alpha'(t)\rangle.\]
\end{lemma}

% Source: Corsin Nick/Class Notes/Week 10.pdf, pp. 4--5
\begin{proposition}[The shortest path is a straight line]
\label{prop:geodesic_straight_line}
Assume that $c \in C^\infty([0,1],\mathbb{R}^n)$ with $c(0)=0$, $c(1)=x$ and
$\lvert c'(t)\rvert \equiv \lambda$ is the optimal path we are looking for. Then $c(t) = xt$
and $\lambda = \lvert x\rvert$.
\end{proposition}

\begin{proof}
Let $g \in C^\infty([0,1],\mathbb{R}^n)$ with $g(0)=g(1)=0$, and define, for $\varepsilon \in
\mathbb{R}$, the family of paths $c_\varepsilon(t) := c(t) + \varepsilon g(t)$, so that
$c_0 = c$. The boundary conditions on $g$ are what keep the whole family admissible:
$c_\varepsilon(0) = c(0) = 0$ and $c_\varepsilon(1) = c(1) = x$ for every $\varepsilon$, so each
$c_\varepsilon$ is a competitor in $\mathcal{A}$ and not merely a nearby curve.

% Sonnet 5 (Medium) --- FIG-W10-02
\begin{center}
\begin{tikzpicture}[>=Latex, scale=1.2]
  \draw[red, thick] (0,0) -- (3.6,1.1) node[midway, below, font=\scriptsize, sloped] {$c$};
  \foreach \k in {1,2,3} {
    \draw[green!45!black, thick, opacity={0.5+0.15*\k}]
      plot[smooth] coordinates {(0,0) (1.1,{0.55+0.13*\k}) (2.3,{0.75+0.10*\k}) (3.6,1.1)};
  }
  \node[font=\scriptsize] at (2.0,1.55) {$c_\varepsilon$};
  \fill (0,0) circle (1.3pt) node[left, font=\scriptsize] {$0$};
  \fill (3.6,1.1) circle (1.3pt) node[right, font=\scriptsize] {$x$};
\end{tikzpicture}
\end{center}

By assumption, $c$ is a minimizer of $L$, so it is a critical point of the family, i.e.
\[\left.\frac{d}{d\varepsilon}\right|_{\varepsilon=0} L(c_\varepsilon) = 0.\]
(This equality itself requires a proof, but the idea is exactly the same as that of
optimization in $\mathbb{R}^n$.) First,
\[\left.\frac{\partial}{\partial\varepsilon}\right|_{\varepsilon=0} c_\varepsilon(t)
  = \left.\frac{\partial}{\partial\varepsilon}\right|_{\varepsilon=0}
    \big(c(t)+\varepsilon g(t)\big) = g(t).\]
Recalling $\lvert c'\rvert \equiv \lambda$ and using
\Cref{lem:inner_product_derivative_identities},
\[
\begin{aligned}
\left.\frac{d}{d\varepsilon}\right|_{\varepsilon=0} L(c_\varepsilon)
  &= \left.\frac{d}{d\varepsilon}\int_0^1 \lvert c_\varepsilon'(t)\rvert\,dt\right|_{\varepsilon=0}
   = \int_0^1 \left.\frac{\partial}{\partial\varepsilon}\lvert c_\varepsilon'(t)\rvert\,dt
     \right|_{\varepsilon=0} \\
  &= \int_0^1 \frac{\big\langle c_\varepsilon'(t),\ \partial_\varepsilon c_\varepsilon'(t)
       \big\rangle}{\lvert c_\varepsilon'(t)\rvert}\,dt\bigg|_{\varepsilon=0}
   = \frac1\lambda \int_0^1 \Big\langle c'(t),\ \partial_\varepsilon\big|_{\varepsilon=0}
       c_\varepsilon'(t)\Big\rangle\,dt \\
  &= \frac1\lambda \int_0^1 \frac{d}{dt}\Big\langle c'(t), g(t)\Big\rangle
       - \big\langle c''(t), g(t)\big\rangle\,dt
   = \frac1\lambda\left[\big\langle c'(t),g(t)\big\rangle\Big|_0^1
       - \int_0^1 \big\langle c''(t),g(t)\big\rangle\,dt\right] \\
  &= -\frac1\lambda \int_0^1 \big\langle c''(t),g(t)\big\rangle\,dt \ \overset{!}{=}\ 0,
\end{aligned}
\]
where the boundary term vanished since $g(0)=g(1)=0$. Since $g$ is arbitrary, it follows that
$c''(t) \equiv 0$. This last step is the \newterm{fundamental lemma of the calculus of
variations}: a continuous function integrating to zero against \emph{every} smooth $g$ vanishing
at the endpoints must itself vanish, since otherwise one could choose $g$ to be a bump
concentrated where $c''$ has a fixed sign and make the integral non-zero. Together with
$c(0)=0$ and $c(1)=x$ this gives $c(t)=xt$ and $\lambda = \lvert x\rvert$.
\end{proof}

% Corsin flags both of these steps in his own proof with a warning triangle but does not resolve
% either. (Was an ainote until 2026-08-09; by Test 1 in style.md it is mathematics, and the
% duplicate \newterm of the fundamental lemma it carried is now left only at its first
% occurrence, inside the proof above.)
\begin{remark}[The two steps that hide real work]
\label{rem:geodesic_hidden_steps}
The proof moves quickly over $\left.\frac{d}{d\varepsilon}\right|_{\varepsilon=0}L(c_\varepsilon)
= 0$ and over the final $\overset{!}{=}\,0$, and those are precisely the two places where
something is being assumed.

The first is the claim that a minimizer of $L$ over the infinite-dimensional set $\mathcal A$ is
a critical point of every one-parameter family through it. That is the same fact used for
optimization on $\mathbb{R}^n$ in \cref{prop:extremum_implies_critical}, and it wants $\mathcal A$
to be a vector space. It is not one, since $\mathcal A$ is an affine subspace. What rescues the
argument is that the \emph{differences} $c_\varepsilon - c$ range over
$\{g \in C^\infty([0,1],\mathbb{R}^n) : g(0)=g(1)=0\}$, which is a genuine vector space, so the
perturbation directions form one even though the curves themselves do not.

The second is the fundamental lemma named in the proof, which upgrades \qt{the integral vanishes
for every $g$} to \qt{$c''$ vanishes at every $t$}, rather than merely on average. Corsin gives
neither its name nor its proof.
\end{remark}

% Generator: Claude Opus 5 (High) --- the hypothesis $\lvert c'\rvert \equiv \lambda$ is used
% twice in the computation (to pull $1/\lambda$ out of the integral) and is never commented on.
\begin{remark}[Why constant speed is assumed, and why it costs nothing]
\label{rem:geodesic_constant_speed}
\Cref{prop:geodesic_straight_line} assumes $\lvert c'(t)\rvert \equiv \lambda$, which looks like
a restriction on the competitor curves and is not one. Length is invariant under
reparametrization, since substituting $t = \varphi(s)$ for an increasing $C^1$ bijection
$\varphi$ leaves $\int\lvert c'\rvert\,dt$ unchanged. Every curve with nowhere-vanishing speed
therefore has a constant-speed representative tracing the same image, obtained by running along
it proportionally to arc length. The hypothesis chooses that representative.

It is worth choosing, because the functional $L$ cannot see the difference and the equation can.
Without it the computation still goes through, but $1/\lvert c'(t)\rvert$ stays inside the
integral and the conclusion is that $\big(c'/\lvert c'\rvert\big)' \equiv 0$, which says the
\emph{direction} of travel is constant while the speed is free. Constant speed converts that
into the cleaner $c'' \equiv 0$. \Cref{ex:ai_variation_without_constant_speed} carries this out.
\end{remark}

\begin{aiexercise}[The fundamental lemma of the calculus of variations]
% Generator: Claude Opus 5 (High)
\label{ex:ai_fundamental_lemma}
The proof of \cref{prop:geodesic_straight_line} passes from \qt{the integral vanishes for every
$g$} to \qt{$c''$ vanishes at every $t$} in one step. Prove that step. Precisely: let
$h \in C^0\big([0,1],\mathbb{R}^n\big)$ satisfy
\[\int_0^1 \big\langle h(t),\,g(t)\big\rangle\,dt = 0
  \qquad\text{for every } g \in C^\infty\big([0,1],\mathbb{R}^n\big) \text{ with } g(0)=g(1)=0.\]
Show that $h \equiv 0$.

\exhint{Argue by contradiction, and use a bump function of the kind constructed in
\cref{def:smooth_bump_function} to concentrate $g$ where $h$ is known not to vanish.}
\end{aiexercise}

\begin{aiexercise}[The variation without the constant-speed hypothesis]
% Generator: Claude Opus 5 (High)
\label{ex:ai_variation_without_constant_speed}
Repeat the computation in the proof of \cref{prop:geodesic_straight_line} for a minimiser
$c \in C^2\big([0,1],\mathbb{R}^n\big)$ that is \emph{not} assumed to have constant speed, only
$\lvert c'(t)\rvert > 0$ for every $t$.
\begin{enumerate}[label=\textbf{(\alph*)}]
  \item Show that $\left.\frac{d}{d\varepsilon}\right|_{\varepsilon=0}L(c_\varepsilon)
    = \int_0^1 \big\langle c'(t)/\lvert c'(t)\rvert,\ g'(t)\big\rangle\,dt$, and integrate by
    parts to bring it into the form $-\int_0^1\langle \cdot,\,g(t)\rangle\,dt$.
  \item Conclude, using \cref{ex:ai_fundamental_lemma}, that
    $\big(c'/\lvert c'\rvert\big)' \equiv 0$.
  \item Deduce that the image of $c$ is still the straight segment from $0$ to $x$, and that
    $L(c) = \lvert x\rvert$. Which feature of $c$ does the argument now fail to pin down, and
    why is that unavoidable?
\end{enumerate}
\end{aiexercise}
